% ===== 原文 pp.76–89（PDF p.91–104）· tr-05 · c02d =====
% 本块承接 c02c（定义 2.36 之后同一句话的后半句 "with $f_{i_0i_2\cdots i_n}$ a holomorphic function on $U$."），
% 按 STYLE_GUIDE §10「从句子中段开始者直接续写」，不重复上文，也不重写下文已由 c02c 写出的内容。
% 覆盖 tag 2.29–2.43（编号公式 2.29–2.43 共 15 条）；本块内无插图。

% [原文 p.76]
其中 $f_{i_0i_2\cdots i_n}$ 是 $U$ 上的全纯函数。设 $\tilde\cX$ 是 $\tilde\cM\times\C^{n+1}$ 中由方程
\[
  e^{t_1}=x_0\cdots x_n
\]
定义的子空间，其中 $x_0,\ldots,x_n$ 是 $\C^{n+1}$ 因子上的坐标。令
\[
  W_0=x_0+\cdots+x_n,
\]
于是可以定义
\begin{equation}
  W=t_0+W_0+\sum_{i=2}^{n}t_iW_0^{i}. \tag{2.29}
\end{equation}
记 $\pi:\tilde\cX\to\tilde\cM$ 为投影。

可以验证，它在 $\tilde\cM$ 的每一点附近确实是万有的；至于 $W$ 的描写方式的具体选取，其细节对我们并不重要。我们特别关心 $\tilde\cM$ 中使 $t_0=t_1=\cdots=t_n=0$ 的点 $0$。

于是可以考虑 $W$ 的临界点构成的子概形 $\operatorname{Crit}(W)\subset\tilde\cX$。它由 $\partial W/\partial x_0,\ldots,\partial W/\partial x_n$ 生成的理想定义。计算可得
\[
  \frac{\partial W}{\partial x_i}
  =\frac{\partial W_0}{\partial x_i}\left(1+\sum_{i=2}^{n}it_iW_0^{i-1}\right).
\]
由于 $t_2,\ldots,t_n$ 是形式参数，$1+\sum_{i=2}^{n}it_iW_0^{i-1}$ 可逆，故由诸 $\partial W/\partial x_i$ 生成的理想与由诸 $\partial W_0/\partial x_i$ 生成的理想相同。于是由例 2.33，$\operatorname{Crit}(W)$ 由轨迹 $x_0=\cdots=x_n=\mu$ 给出，其中 $\mu$ 取遍 $e^{t_1}$ 的 $(n+1)$ 次根。因此 $\operatorname{Crit}(W)$ 在 $\tilde\cM$ 上是 étale 的，且 $\pi_*\cO_{\operatorname{Crit}(W)}$ 是秩为 $n+1$ 的局部自由 $\cO_{\tilde\cM}$-模层。在 $\tilde\cM$ 上可以找到截面 $p_1,\ldots,p_{n+1}:\tilde\cM\to\operatorname{Crit}(W)$。

% [Tutor Note: 这里“$\operatorname{Crit}(W)$ 在 $\tilde\cM$ 上 étale、$\pi_*\cO_{\operatorname{Crit}(W)}$ 秩为 $n+1$”必须按形式意义理解：$\tilde\cM$ 是形式环化空间（定义 2.36），只有 $t_1$ 是真正的坐标，$t_2,\ldots,t_n$ 是形式参数，故“$n+1$ 个互不相同的根”指 $x^{n+1}=e^{t_1}$ 在 $\C[[t_1]]$ 中的 $n+1$ 个根（$e^{t_1}$ 为单位，各根互异）。核对锚点：$x_0=\cdots=x_n=\mu$ 且 $\mu^{n+1}=e^{t_1}$。]

在 $\tilde\cM\times\C^\times$ 上，以 $\hbar$ 记 $\C^\times$ 的坐标，我们如下定义一个 $\C$-向量空间的局部系统 $R$。注意，局部系统只涉及 $\tilde\cM\times\C^\times$ 的底层拓扑空间，而与结构层 $\cO_{\tilde\cM}$ 无关。明白了这一点，$R$ 在 $(t_1,\hbar)\in\C\times\C^\times$ 上的纤维将是
\[
  H_n\bigl(\pi^{-1}(t_1),\Re(W|_{\pi^{-1}(t_1)}/\hbar)\ll0;\C\bigr).
\]
这里 $W|_{\pi^{-1}(t_1)}=x_0+\cdots+x_n$，其中 $x_0\cdots x_n=e^{t_1}$，因为我们忽略了变量 $t_0,t_2,\ldots,t_n$。该函数有 $n+1$ 个临界点。因此不必担心该空间的维数发生跳跃。

令 $\cR=R\otimes_\C\cO_{\tilde\cM\times\C^\times}$；它是 $\tilde\cM\times\C^\times$ 上秩为 $n+1$ 的局部自由层，并带有 Gauss--Manin 联络 $\nabla^{GM}$，其平坦截面正是 $R$ 的截面。于是还有对偶的局部自由层
\[
  \cR^\vee=\mathcal{H}om_{\cO_{\tilde\cM\times\C^\times}}(\cR,\cO_{\tilde\cM\times\C^\times}),
\]
其对偶局部系统 $R^\vee\subseteq\cR^\vee$ 同样由某个联络 $\nabla^{GM}$ 的平坦截面给出。

注意，$\cR^\vee$ 在开集 $U\times V$（其中 $U\subseteq\tilde\cM$、$V\subseteq\C^\times$ 为开集）上的截面可以用微分形式显式地给出：
\[
  f\Omega,
\]

% [原文 p.77]
其中 $f$ 是 $\pi^{-1}(U)\times V$ 上的全纯函数，且对每个 $u\in U$、$v\in V$，$f|_{\pi^{-1}(u)\times\{v\}}$ 是代数的，并且
\[
  \Omega=\frac{dx_1\wedge\cdots\wedge dx_n}{x_1\cdots x_n}.
\]
形式 $f\Omega$ 在 $U\times V$ 上定义了映射 $\cR\to\cO_{\tilde\cM\times\C^\times}$：把 $\cR$ 在 $U\times V$ 上的截面 $\Xi$ 映为
\[
  \int_\Xi e^{W/\hbar}f\Omega\in\cO_{\tilde\cM\times\C^\times}.
\]
我们用 $[f\Omega]$ 记以这种方式确定的 $\cR^\vee$ 的截面。

下面把这件事在我们所处的形式设定中的含义说得更精确些。利用 (2.29) 中 $W$ 的表达式，写
\[
  e^{W/\hbar}=e^{(x_0+\cdots+x_n)/\hbar}g,
\]
其中 $g$ 是变量 $t_0,t_2,\ldots,t_n$ 的形式幂级数，其中出现的每个系数都是 $x_0,\ldots,x_n$ 与 $\hbar$ 的代数函数。于是该展开中每一项的积分都收敛，即等于
\[
  \int_\Xi e^{(x_0+\cdots+x_n)/\hbar}h\Omega,
\]
其中 $h$ 对每个 $u\in U$、$v\in V$ 都满足 $h|_{\pi^{-1}(u)\times\{v\}}$ 是代数的。事实上，$e^{(x_0+\cdots+x_n)/\hbar}$ 的指数衰减控制住 $h$ 的多项式增长。于是可见 $f\Omega$ 定义了 $\cR^\vee$ 的一个元素。

有一个标准的技巧可以在形式层面描述 Gauss--Manin 联络，这里略作勾勒。给定 $\pi^{-1}(U)\times V$ 上的 $n$-形式 $f\Omega$，$f\Omega$ 表示 $\cR^\vee$ 的一个截面。现在把全空间 $\tilde\cX\times\C^\times$ 看作带有 Landau--Ginzburg 势 $\hbar^{-1}W$ 的空间。若 $f\Omega$ 在 $\pi^{-1}(U)\times V$ 上关于 $d+d(\hbar^{-1}W)\wedge$ 是闭的，那么 $f\Omega$ 实际上表示 $\cR^\vee$ 的一个平坦截面。Gauss--Manin 联络通过把 $(n+1)$-形式 $(d+d(\hbar^{-1}W)\wedge)(f\Omega)$ 作用到由 $\tilde\cM\times\C^\times$ 提升而来的 $\tilde\cX\times\C^\times$ 上的切向量 $X$，来度量它在全空间上不为零的程度。特别地，若 $X$ 实际上是某个切向量从 $\tilde\cM$ 到 $\tilde\cX\times\C^\times$ 的提升，且在 $\C^\times$ 方向上的投影为零，则
\[
  \iota(X)(d+d(\hbar^{-1}W)\wedge)(f\Omega)
  =\bigl(X(f)+\hbar^{-1}X(W)f\bigr)\Omega,
\]
而
\[
  \iota(\hbar\partial_\hbar)(d+d(\hbar^{-1}W)\wedge)(f\Omega)
  =(\hbar\partial_\hbar f-\hbar^{-1}Wf)\Omega.
\]
由此得到
\begin{equation}
  \nabla^{GM}_X[f\Omega]=\bigl[(X(f)+\hbar^{-1}X(W)f)\Omega\bigr] \tag{2.30}
\end{equation}
以及
\begin{equation}
  \nabla^{GM}_{\hbar\partial_\hbar}[f\Omega]
  =\bigl[(\hbar\partial_\hbar f-\hbar^{-1}Wf)\Omega\bigr]. \tag{2.31}
\end{equation}
容易验证，这些类只依赖于 $f\Omega$ 的类，而不依赖于 $f\Omega$ 本身或提升 $X$。

% [Tutor Note: 核对 (2.30)、(2.31) 的符号约定：$\iota(X)$ 与 $d(\hbar^{-1}W)\wedge$ 缩并给出 $\hbar^{-1}X(W)f\Omega$ 一项；而 $\iota(\hbar\partial_\hbar)(d(\hbar^{-1}W)\wedge f\Omega)=\hbar\partial_\hbar(\hbar^{-1}W)f\Omega=-\hbar^{-1}Wf\Omega$，其中用 $\hbar\partial_\hbar(\hbar^{-1})=-\hbar^{-1}$，这就是 (2.31) 中负号的来源。两式的相容性检验：把 $f$ 换成 $\hbar f$（$\hbar$ 是 $\cO\{\hbar\}$ 的系数，故 $\nabla^{GM}$ 可穿透），代入 (2.30)、(2.31) 应得到同一结果。]

\begin{definition}
我们如下定义 $\tilde\cM$ 上的 B-模型半无限 Hodge 结构变分。丛 $\cE$ 是 $\cO_{\tilde\cM}\{\hbar\}$-模，使得 $\Gamma(U,\cE)$ 由 $\cR^\vee$ 在形如
\[
  U\times\{\hbar\mid|\hbar|<\epsilon\}
\]
的开集上的截面组成，

% [原文 p.78]
这些截面由形式 $f\Omega$ 给出，其中 $f$ 在 $|\hbar|<\epsilon$（包括 $\hbar=0$）上全纯，且限制到 $\pi^{-1}(u)\times\{\hbar\}$ 上时代数。联络
\[
  \nabla:\cE\to\Omega^1_{\tilde\cM}\otimes\hbar^{-1}\cE
\]
由 $\cR^\vee$ 上的 Gauss--Manin 联络给出，即对 $\cE$ 的截面 $s$，
\[
  \nabla_Xs=\nabla^{GM}_Xs,
\]
其中右端的 $X$ 是 $X$ 到 $\tilde\cM\times\C^\times$ 的提升，且在 $\C^\times$ 方向上的投影为零。
\end{definition}

另一种理解方式如下：我们可以把局部自由层 $\cR^\vee$ 从 $\tilde\cM\times\C^\times$ 延拓到 $\tilde\cM\times\C$，规定 $\cR^\vee$ 的截面 $[f\Omega]$ 可延拓当且仅当 $f$ 能作为全纯函数穿过 $\hbar=0$ 延拓。两个形式 $f_1\Omega$ 与 $f_2\Omega$ 确定同一截面，当且仅当对 $0\in\C$ 的邻域中的每个 $\hbar$，$(f_1-f_2)\Omega$ 落在 $\hbar d+dW\wedge$ 的像中；注意当 $\hbar\ne0$ 时，这与落在 $d+\hbar^{-1}dW\wedge$ 的像中是一回事。由 [102] 可知，$\cR^\vee$ 的这个延拓实际上是一个向量丛。作为向量丛，$\cR^\vee$ 在 $(u,0)$ 上的纤维是 $\mathbb{H}^n(\pi^{-1}(u),(\Omega^{n}_{\pi^{-1}(u)},dW\wedge))$。该向量丛在 $\hbar=0$ 附近的截面给出 $\cE$ 的截面。特别地，$\cE/\hbar\cE$ 是秩为 $n+1$ 的向量丛，其纤维正是 $\mathbb{H}^n(\pi^{-1}(u),(\Omega^{n}_{\pi^{-1}(u)},dW\wedge))$。

令
\[
  (-):\tilde\cM\times\C^\times\to\tilde\cM\times\C^\times
\]
记映射
\[
  (u,\hbar)\mapsto(u,-\hbar).
\]
则配对 (2.25) 诱导出配对
\[
  (\cdot,\cdot):(-)^*\cR^\vee\times\cR^\vee\to\cO_{\tilde\cM\times\C^\times}.
\]
由此得到 $\cE$ 上的配对 $(\cdot,\cdot)_\cE$，定义为
\[
  (s_1,s_2)_\cE(\hbar)
  =\frac{(-1)^{n(n+1)/2}}{(2\pi i\hbar)^n}\bigl((-)^*s_1,s_2\bigr).
\]

至于分次算子，我们先取 $\tilde\cM$ 上的向量场 $E$，它由
\[
  E=(n+1)\partial_{t_1}+\sum_{i=0}^{n}(1-i)t_i\partial_{t_i}
\]
给出。然后如下定义分次算子：
\begin{equation}
  \operatorname{Gr}(s)=\nabla^{GM}_{\hbar\partial_\hbar+E}(s)-s. \tag{2.32}
\end{equation}

\begin{proposition}
数据 $\cE$、$\nabla$ 与 $(\cdot,\cdot)_\cE$ 给出一个半无限 Hodge 结构变分。算子 $\operatorname{Gr}$ 是分次算子。
\end{proposition}

\para{证明。}首先注意，由 (2.30)，$\nabla$ 确实把 $\cE$ 映入 $\Omega^1_{\tilde\cM}\otimes\hbar^{-1}\cE$。其次，利用 Lefschetz 针（Lefschetz thimble）的基 $\{\Delta_p^\pm\}$ 可以描述配对 $(\cdot,\cdot)_\cE$。由 (2.26) 可写
\begin{equation}
  (s_1,s_2)_\cE
  =\frac{(-1)^{n(n+1)/2}}{(2\pi i\hbar)^n}\sum_p
  \left(\int_{\Delta^-_p}s_1(-\hbar)\right)
  \left(\int_{\Delta^+_p}s_2(\hbar)\right). \tag{2.33}
\end{equation}

% [原文 p.79]
于是由驻相近似（命题 2.35）得到
\begin{align*}
  ([f\Omega],[g\Omega])_\cE
  &=\pm\sum_p\frac{f(p,0)g(p,0)}{\operatorname{Hess}(W(p))}+O(\hbar)\\
  &\in\cO_{\tilde\cM}\{\hbar\}.
\end{align*}
所以 $(\cdot,\cdot)_\cE$ 取值于 $\cO_{\tilde\cM}\{\hbar\}$，这正是所需的。

由于 Gauss--Manin 联络按定义是平坦的，$\nabla$ 是平坦的。定义 2.20 的公理 (1)--(3) 显然由定义满足。为看出 (4)，我们如下进行。回忆有 $\operatorname{Crit}(W)\subset\tilde\cX$，以及投影 $\pi:\operatorname{Crit}(W)\to\tilde\cM$ 的截面 $p:\tilde\cM\to\operatorname{Crit}(W)$，它们给出一族临界点。这样的截面共有 $n+1$ 个。对每个 $p$，取 $\tilde\cX$ 上的全纯函数 $f_p$，它在 $\pi$ 的纤维上正则，且满足 $f_p\circ p'=\delta_{pp'}$，于是作为 $\operatorname{Crit}(W)$ 上的函数，$\{f_p\}$ 构成向量丛 $\pi_*\cO_{\operatorname{Crit}(W)}$ 的截面的一组基。于是
\begin{align*}
  ([f_{p'}\Omega],[f_{p''}\Omega])_\cE
  &=\pm\sum_p\frac{f_{p'}(p,0)f_{p''}(p,0)}{\operatorname{Hess}(W(p))}+O(\hbar)\\
  &=\pm\frac{\delta_{p'p''}}{\operatorname{Hess}W(p')}+O(\hbar).
\end{align*}
现在由构造，对每个点 $u\in\tilde\cM$，形式 $f_p\Omega$ 构成
\[
  \mathbb{H}^n\bigl(\pi^{-1}(u),(\Omega^{\bullet}_{\pi^{-1}(u)},dW\wedge)\bigr)
\]
的一组基，因为函数 $f_p$ 显然构成 $W|_{\pi^{-1}(u)}$ 的 Milnor 环的一组基。因此类 $[f_p\Omega]$ 也给出秩为 $n+1$ 的向量丛 $\cE/\hbar\cE$ 的一个局部基，因为该向量丛的纤维正是
\[
  \mathbb{H}^n\bigl(\pi^{-1}(u),(\Omega^{\bullet}_{\pi^{-1}(u)},dW\wedge)\bigr).
\]
于是定义 2.20 的非退化条件 (4) 成立，这一点是清楚的。

% [Tutor Note: 从 (2.33) 到本页第一式的 $\pm$ 来自驻相近似（命题 2.35）中 $\hbar\to0$ 的角度扇形选取与 $\sqrt{\operatorname{Hess}W(p)}$ 的分支；全段应以同一个 $\pm$ 整体解读。核对锚点：第二式把分子在 $p=p'=p''$ 处取 $\delta_{p'p''}$，分母保持 $\operatorname{Hess}W(p')$，$\pm$ 并未消去，因此 (2.40) 中 $i+j\le n$ 时的非零常数只确定到符号。]

下面转向分次算子。首先注意
\begin{equation}
  \pi_*\left(\sum_{i=0}^{n}x_i\partial_{x_i}\right)=(n+1)\partial_{t_1}, \tag{2.34}
\end{equation}
其中用到 $e^{t_1}=x_0\cdots x_n$。于是可以把 $E$ 提升为 $\tilde\cX$ 上的向量场
\[
  \sum_{i=0}^{n}x_i\partial_{x_i}+\sum_{i=0}^{n}(1-i)t_i\partial_{t_i},
\]
并计算
\begin{align*}
  E(W)&=\left(\sum_{i=0}^{n}x_i\partial_{x_i}+\sum_{i=0}^{n}(1-i)t_i\partial_{t_i}\right)
  \left(W_0+\sum_{\substack{i=0\\ i\ne1}}^{n}t_iW_0^{i}\right)\\
  &=W_0+\sum_{\substack{i=0\\ i\ne1}}^{n}\bigl(i+(1-i)\bigr)t_iW_0^{i}\\
  &=W.
\end{align*}
事实上容易看出，$E$ 是 $\tilde\cM$ 上唯一能提升为保持 $W$ 的 $\tilde\cX$ 上向量场的向量场。

% [原文 p.80]
把 $\operatorname{Gr}$ 作用于 $\cE$ 的截面 $[f\Omega]$，利用 (2.30)、(2.31) 以及 $E(W)=W$，可见
\[
  \operatorname{Gr}\bigl([f\Omega]\bigr)
  =\bigl[(\hbar\partial_\hbar f+E(f)-f)\Omega\bigr].
\]
于是确实有 $\operatorname{Gr}(\cE)\subseteq\cE$。现在定义 2.20 的公理 (5) 是显然的，(6) 由 $\nabla^{GM}$ 的平坦性得到，而 (7) 的验证如下。由 (2.33)，有
\begin{align*}
  (\hbar\partial_\hbar+E)(s_1,s_2)_\cE
  &=(\hbar\partial_\hbar+E)\left(\frac{(-1)^{n(n+1)/2}}{(2\pi i\hbar)^n}
  \sum_p\left(\int_{\Delta^-_p}s_1(-\hbar)\right)\left(\int_{\Delta^+_p}s_2(\hbar)\right)\right)\\
  &=-n(s_1,s_2)_\cE\\
  &\quad+\frac{(-1)^{n(n+1)/2}}{(2\pi i\hbar)^n}
  \left(\int_{\Delta^-_p}(\nabla^{GM}_{\hbar\partial_\hbar+E}s_1)(-\hbar)\right)
  \left(\int_{\Delta^+_p}s_2(\hbar)\right)\\
  &\quad+\frac{(-1)^{n(n+1)/2}}{(2\pi i\hbar)^n}
  \left(\int_{\Delta^-_p}s_1(-\hbar)\right)
  \left(\int_{\Delta^+_p}\nabla^{GM}_{\hbar\partial_\hbar+E}s_2(\hbar)\right)\\
  &=-n(s_1,s_2)_\cE\\
  &\quad+\frac{(-1)^{n(n+1)/2}}{(2\pi i\hbar)^n}
  \left(\int_{\Delta^-_p}\bigl(\operatorname{Gr}(s_1)+s_1\bigr)(-\hbar)\right)
  \left(\int_{\Delta^+_p}s_2(\hbar)\right)\\
  &\quad+\frac{(-1)^{n(n+1)/2}}{(2\pi i\hbar)^n}
  \left(\int_{\Delta^-_p}s_1(-\hbar)\right)
  \left(\int_{\Delta^+_p}\bigl(\operatorname{Gr}(s_2)+s_2\bigr)(\hbar)\right)\\
  &=(\operatorname{Gr}(s_1),s_2)_\cE+(s_1,\operatorname{Gr}(s_2))_\cE+(2-n)(s_1,s_2)_\cE.
\end{align*}
于是 $\operatorname{Gr}$ 是分次算子，且 $D=2-\dim_\C X$。

\subsubsection{B-模型的 Frobenius 流形}现在我们已经有了上一小节的 $\tilde\cM$ 上的 B-模型半无限 Hodge 结构变分 $\cE$、$\nabla$、$(\cdot,\cdot)_\cE$ 以及分次算子 $\operatorname{Gr}$。按 §2.1.7，再加上对某些数据的适当额外选取，便会在 $\tilde\cM$ 上得到 Frobenius 流形结构。为了使它与由 $X=\PP^n$ 得到的 A-模型 Frobenius 流形一致，这些数据不能任意选取。然而，为了具体给出这些数据，我们需要更细致地考察局部系统 $R$ 的性质。特别地，我们希望理解该局部系统的单值性。由于 $R$ 是 $\tilde\cM\times\C^\times$ 上的局部系统，而 $\tilde\cM$ 是单连通的，唯一有趣的单值性来自 $\C^\times$ 因子中的回路，因此事实上可以限制到点 $0\in\tilde\cM$，考虑 $\C^\times$ 上纤维为 $H_n(\check\cX,\Re W_0/\hbar\ll0;\C)$ 的局部系统，其中 $\check\cX=V(x_0\cdots x_n-1)$，$W_0=x_0+\cdots+x_n$。

\begin{proposition}
设 $\check\cX=V(x_0\cdots x_n-1)\subseteq\operatorname{Spec}\C[x_0,\ldots,x_n]$，$W_0=x_0+\cdots+x_n$。则对 $\C^\times$ 上局部系统 $R$ 的任一局部截面 $\Xi$（其在 $\hbar\in\C^\times$ 上的纤维由 $H_n(\check\cX,\Re W_0/\hbar\ll0;\C)$ 给出），积分
\[
  \psi=\int_\Xi e^{W_0/\hbar}\Omega
\]
满足常微分方程
\begin{equation}
  \left(-\frac{1}{n+1}\hbar\partial_\hbar\right)^{n+1}\psi=\hbar^{-(n+1)}\psi. \tag{2.35}
\end{equation}
\end{proposition}

% [原文 p.81]
\para{证明。}令 $\omega_i=x_0\cdots x_{i-1}\Omega$。令 $\Omega_j=\iota(x_j\partial_{x_j})\Omega$。把 $d+\hbar^{-1}dW_0\wedge$ 作用于 $x_0\cdots x_{i-1}\Omega_j$，记住 $x_0\cdots x_n=1$，便得到 $n$-形式
\[
  \bigl((x_j\partial_{x_j}-x_0\partial_{x_0})(x_0\cdots x_{i-1})
  +\hbar^{-1}(x_j-x_0)(x_0\cdots x_{i-1})\bigr)\Omega,
\]
特别地，
\[
  \bigl((x_j\partial_{x_j}-x_i\partial_{x_i})(x_0\cdots x_{i-1})
  +\hbar^{-1}(x_j-x_i)(x_0\cdots x_{i-1})\bigr)\Omega
\]
落在 $d+\hbar^{-1}dW_0\wedge$ 的像中。于是在 $\mathbb{H}^n(\check\cX,(\Omega^{\bullet}_{\check\cX},d+\hbar^{-1}dW_0\wedge))$ 中，有
\[
  \bigl[(x_j\partial_{x_j})(x_0\cdots x_{i-1})\Omega\bigr]
  =-\bigl[\hbar^{-1}(x_j-x_i)\omega_i\bigr].
\]
对所有的 $j$（$0\le j\le n$）求和，得到等式
\[
  [i\omega_i]=-\bigl[\hbar^{-1}(W_0-(n+1)x_i)\omega_i\bigr].
\]
由于由 (2.31) 有 $\nabla^{GM}_{\hbar\partial_\hbar}[\omega_i]=-[\hbar^{-1}W_0\omega_i]$，我们得到
\[
  \frac{1}{n+1}\bigl(-\nabla^{GM}_{\hbar\partial_\hbar}+i\bigr)[\omega_i]
  =\hbar^{-1}[\omega_{i+1}].
\]
由此归纳可得
\[
  \frac{1}{(n+1)^{i}}\bigl(-\nabla^{GM}_{\hbar\partial_\hbar}\bigr)^{i}[\omega_0]
  =\hbar^{-i}[\omega_i].
\]
由于 $x_0\cdots x_n=1$ 给出 $\omega_{n+1}=\omega_0$，我们得到
\[
  \frac{1}{(n+1)^{n+1}}\bigl(-\nabla^{GM}_{\hbar\partial_\hbar}\bigr)^{n+1}[\omega_0]
  =\hbar^{-(n+1)}[\omega_0].
\]
于是 $\omega_0$ 在 $\Xi$ 上的积分满足方程 (2.35)。

注意 (2.35) 是 $(n+1)$ 阶常微分方程，故有 $n+1$ 个线性无关的解。我们可以相当容易地写出一个基本解组。

\begin{proposition}
在环 $\C[\alpha]/(\alpha^{n+1})$ 上工作，写
\[
  \hbar^{-(n+1)\alpha}
  =\exp\bigl(-(n+1)\alpha\log\hbar\bigr)
  =\sum_{i=0}^{n}\frac{\bigl(-(n+1)\bigr)^{i}\log^{i}\hbar}{i!}\alpha^{i}.
\]
则表达式
\begin{equation}
  \xi(\hbar,\alpha)
  =\hbar^{-(n+1)\alpha}\sum_{d=0}^{\infty}\hbar^{-(n+1)d}
  \prod_{i=1}^{d}\frac{1}{(\alpha+i)^{n+1}} \tag{2.36}
\end{equation}
中 $1,\alpha,\ldots,\alpha^{n}$ 的系数构成 (2.35) 的一个基本解组。
\end{proposition}

\para{证明。}系数 $\alpha^{i}$（$i=0,\ldots,n$）是线性无关的，这一点可以这样看出：$\alpha^{i}$ 的系数形如 $f\log^{i}\hbar$ 加上 $\log\hbar$ 的低阶项，其中 $f\ne0$。为看出这些系数确实是解，注意对 $\C^\times$ 上取值于 $\C[\alpha]/(\alpha^{n+1})$ 的函数 $\varphi$，有
\[
  -\frac{1}{n+1}\hbar\partial_\hbar
  \bigl(\hbar^{-(n+1)\alpha}\varphi\bigr)
  =\alpha\hbar^{-(n+1)\alpha}\varphi
  -\frac{1}{n+1}\hbar^{-(n+1)\alpha}\hbar\partial_\hbar\varphi.
\]
于是，若 $\varphi$ 满足方程
\[
  \left(\alpha-\frac{1}{n+1}\hbar\partial_\hbar\right)^{n+1}\varphi
  =\hbar^{-(n+1)}\varphi,
\]
则 $\psi=\hbar^{-(n+1)\alpha}\varphi$ 满足 (2.35)。

% [原文 p.82]
取
\[
  \varphi=\sum_{d=0}^{\infty}\hbar^{-(n+1)d}
  \prod_{i=1}^{d}\frac{1}{(\alpha+i)^{n+1}},
\]
可见
\[
  \left(\alpha-\frac{1}{n+1}\hbar\partial_\hbar\right)\varphi
  =\sum_{d=0}^{\infty}(\alpha+d)\hbar^{-(n+1)d}
  \prod_{i=1}^{d}\frac{1}{(\alpha+i)^{n+1}},
\]
从而
\[
  \left(\alpha-\frac{1}{n+1}\hbar\partial_\hbar\right)^{n+1}\varphi
  =\sum_{d=1}^{\infty}\hbar^{-(n+1)d}
  \prod_{i=1}^{d-1}\frac{1}{(\alpha+i)^{n+1}},
\]
（注意对 $d=0$，$(\alpha+d)^{n+1}=0$）而
\[
  \hbar^{-(n+1)}\varphi
  =\sum_{d=0}^{\infty}\hbar^{-(n+1)(d+1)}
  \prod_{i=1}^{d}\frac{1}{(\alpha+i)^{n+1}}.
\]
比较 $\hbar^{-(n+1)d}$ 的系数，两者一致。

% [Tutor Note: (2.35) 左端带负号，正是 (2.31) 中 $\nabla^{GM}_{\hbar\partial_\hbar}$ 的符号约定所致：由 (2.31)，$-\nabla^{GM}_{\hbar\partial_\hbar}[\omega]=[\hbar^{-1}W_0\omega]$，故 $-\frac{1}{n+1}\hbar\partial_\hbar$ 与 $\alpha$ 配对出现（见本页末的 $\left(\alpha-\frac{1}{n+1}\hbar\partial_\hbar\right)^{n+1}\varphi=\hbar^{-(n+1)}\varphi$）。核对锚点：把 (2.36) 的 $\xi$ 写成 $\hbar^{-(n+1)\alpha}\varphi$，代入 (2.35) 应化为关于 $\varphi$ 的同一方程；阶数为 $n+1$ 也与 $\alpha^{n+1}=0$ 的截断一致。]

由该命题可得如下推论：

\begin{lemma}
若 $\Xi_0,\ldots,\Xi_n$ 是 $R$ 的截面的一个局部基，则积分 $\int_{\Xi_i}e^{W_0/\hbar}\Omega$（$i=0,\ldots,n$）构成 (2.35) 的一个基本解组。
\end{lemma}

\para{证明。}回忆定义 2.37 中描述的 $\cR^\vee$ 跨 $\hbar=0$ 的延拓。由于 $1,W_0,\ldots,W_0^{n}$ 张成 $W_0$ 的 Jacobian 环，由类 $[W_0^{i}\Omega]$（$0\le i\le n$）给出的向量丛 $\cR^\vee|_{\{0\}\times\C}$（其中 $0\in\tilde\cM$ 是原点）的截面给出 $\cR^\vee|_{\{0\}\times\C}$ 在 $\hbar=0$ 处纤维的一组基，从而在某个邻域 $V=\{\hbar\in\C\mid|\hbar|<\epsilon\}$ 内给出 $\cR^\vee|_{\{0\}\times\C}$ 的一组截面基。特别地，由于 (2.24) 是完美配对，$(n+1)\times(n+1)$ 矩阵
\[
  \left(\int_{\Xi_i}W_0^{j}e^{W_0/\hbar}\right)_{0\le i,j\le n}
\]
在 $V$ 的每一点处秩为 $n+1$。又
\begin{equation}
  \left(-\hbar^{2}\frac{\partial}{\partial\hbar}\right)^{j}\int_{\Xi_i}e^{W_0/\hbar}\Omega
  =\int_{\Xi_i}W_0^{j}e^{W_0/\hbar}\Omega. \tag{2.37}
\end{equation}
于是 (2.35) 的解组
\[
  \int_{\Xi_0}e^{W_0/\hbar}\Omega,\ \ldots,\ \int_{\Xi_n}e^{W_0/\hbar}\Omega
\]
的 Wronskian 行列式在 $V$ 上不为零，故这些积分给出一个基本解组。

特别地，可以选取 $R$ 的一个局部基 $\Xi_0,\ldots,\Xi_n$，使得 $\int_{\Xi_i}e^{W_0/\hbar}\Omega$ 恰好与 $\xi$ 中 $\alpha^{i}$ 的系数一致，即
\begin{equation}
  \sum_{i=0}^{n}\alpha^{i}\int_{\Xi_i}e^{W_0/\hbar}\Omega
  =\hbar^{-(n+1)\alpha}\sum_{d=0}^{\infty}\hbar^{-(n+1)d}
  \prod_{i=1}^{d}\frac{1}{(\alpha+i)^{n+1}}
  \mod\alpha^{n+1}. \tag{2.38}
\end{equation}
今后我们用 $\Xi_0,\ldots,\Xi_n$ 记 $R$ 的这个特定的基。

% [原文 p.83]
\begin{remark}
我们已经见过几种选取循环基的方式。一般地，我们有循环 $\Delta_p^{+}$，其中 $p$ 取遍 $W_0$ 的临界点。这组基与基 $\Xi_0,\ldots,\Xi_n$ 几乎没有关系。另一方面，在二维情形，我们在例 2.34 中显式描述了 $H_2(\check\cX,\Re W_0/\hbar\ll0;\C)$ 的一组基。用 $\Xi_0',\Xi_1',\Xi_2'$ 记那里描述的基，事实上这两组基之间的关系可以通过直接计算积分来理解。对 $\Xi_i'$ 显式地计算那些振荡积分是一件颇有乐趣的练习。可得如下关系：
\[
  \begin{aligned}
    \Xi_0&=\frac{1}{(2\pi i)^{2}}\Xi_0'\\
    \Xi_1&=\frac{1}{2\pi i}\Xi_1'+\frac{3\gamma}{(2\pi i)^{2}}\Xi_0'\\
    \Xi_2&=\Xi_2'+\frac{3\gamma}{2\pi i}\Xi_1'
    -\frac{1}{(2\pi i)^{2}}\left(\frac{\pi^{2}}{4}+\frac{27}{2}\gamma^{2}\right)\Xi_0'.
  \end{aligned}
\]
这里
\[
  \gamma=\lim_{n\to\infty}\left(-\log(n)+\sum_{k=1}^{n}\frac{1}{k}\right)
\]
是 Euler 常数。$\Xi_0'$ 上的积分是平凡的，来自一次留数计算。$\Xi_1'$ 上的积分要复杂一些，涉及指数积分；而第三个积分是真正的挑战。不过，可以应用 Iritani [61] 的结果，用 $\PP^{2}$ 的结构层的某些特征类来刻画这第三个积分。
\end{remark}

% [Tutor Note: 此处 $\gamma$ 是 Euler--Mascheroni 常数（$\approx0.5772$），不是 Euler 示性数；注意三式构成的矩阵是上三角的，对角元依次为 $(2\pi i)^{-2}$、$(2\pi i)^{-1}$、$1$。核对锚点：p.84 起的单值化论证要求 $\Xi_k$ 与 $\log^{k}\hbar$ 的幂次一一对应，故 $3\gamma$ 只可能出现在严格上三角的位置。]

现在 $\Xi_0,\ldots,\Xi_n$ 被定义为 $\C^\times$ 上局部系统 $R$ 的局部截面。当我们在 $\hbar$-平面中沿逆时针回路行进并跟踪截面 $\Xi_0,\ldots,\Xi_n$，回到回路的起点时，会发现诸 $\Xi_i$ 已变换为诸 $\Xi_i$ 的某个线性组合。这个线性变换就是局部系统 $R$ 的单值变换。为确切看出这个变换是什么，可以研究积分
\[
  \xi_i=\int_{\Xi_i}e^{W_0/\hbar}\Omega
\]
的行为。由于 $e^{W_0/\hbar}\Omega$ 当然给出 $\cR^\vee$ 的一个单值截面，积分的多值性必定恰好来自诸 $\Xi_i$ 的单值变换。

这种多值性如下。注意 $\xi=\sum_{i=0}^{n}\alpha^{i}\xi_i=\hbar^{-(n+1)\alpha}\varphi$，其中 $\varphi$ 是某个单值函数。当我们在 $\C^\times$ 平面中沿绕原点的逆时针回路行进时，$\hbar^{-(n+1)\alpha}=\exp\bigl(-(n+1)\alpha\log\hbar\bigr)$ 被替换为
\[
  \exp\bigl(-(n+1)\alpha(\log\hbar+2\pi i)\bigr)
  =\hbar^{-(n+1)\alpha}\exp\bigl(-(n+1)2\pi i\alpha\bigr).
\]
由此可见，绕该回路一周后，$\xi$ 被替换为 $\xi e^{-(n+1)2\pi i\alpha}$，而 $\xi_k$ 的变换为
\[
  \xi_k\mapsto\sum_{j=0}^{k}
  \frac{\bigl(-(n+1)2\pi i\bigr)^{k-j}}{(k-j)!}\xi_j.
\]
这意味着循环 $\Xi_k$ 的变换方式完全相同：
\[
  \Xi_k\mapsto\sum_{j=0}^{k}
  \frac{\bigl(-(n+1)2\pi i\bigr)^{k-j}}{(k-j)!}\Xi_j.
\]

% [原文 p.84]
为描述 $R^\vee$ 上的对偶单值性，方便的做法是把 $\Xi_0,\ldots,\Xi_n$ 的对偶基写成 $1=\alpha^{0},\alpha^{1},\ldots,\alpha^{n}\in\C[\alpha]/(\alpha^{n+1})$。这使我们能够例如用这组基把 $\cR^\vee$ 的截面 $[f\Omega]$ 写成
\[
  [f\Omega]=\sum_{i=0}^{n}\alpha^{i}\int_{\Xi_i}fe^{W/\hbar}\Omega.
\]
于是单值性在这组对偶基上的作用由转置逆变换给出，容易看出它是
\[
  \alpha^{k}\mapsto\exp\bigl((n+1)2\pi i\alpha\bigr)\alpha^{k}.
\]
特别地，$\hbar^{-(n+1)\alpha}\alpha^{k}$（$k=0,\ldots,n$）是 $\cR^\vee$ 的单值、整体定义的截面，因为 $\alpha^{k}$ 的多值性被 $\hbar^{-(n+1)\alpha}$ 的多值性抵消了。

现在过渡到整个模空间 $\tilde\cM$。局部截面 $\Xi_0,\ldots,\Xi_n$ 现在延拓为整个 $\tilde\cM\times\C^\times$ 上局部系统 $R$ 的多值截面，对偶基 $1,\alpha,\ldots,\alpha^{n}$ 亦然，从而给出 $\cR^\vee$ 在 $\tilde\cM\times\C^\times$ 上整体定义的截面 $\hbar^{-(n+1)\alpha}\alpha^{k}$。于是可以把 $\hbar^{-(n+1)\alpha}\alpha^{k}$ 看作 $\cE\otimes_{\cO_{\tilde\cM}\{\hbar\}}\cO_{\tilde\cM}\{\hbar,\hbar^{-1}\}$ 的截面。此外，由于 $\alpha^{k}$ 是平坦的，$\nabla(\hbar^{-(n+1)\alpha}\alpha^{k})=0$，因此
\[
  \bigl\{\hbar^{-(n+1)\alpha}\alpha^{i}\mid i=0,\ldots,n\bigr\}
\]
构成 $\cH$ 作为自由 $\C\{\hbar,\hbar^{-1}\}$-模的一组基。

现在可以精确地说出我们将用来确定 $\tilde\cM$ 上 Frobenius 流形结构的数据：
\begin{itemize}[leftmargin=1.6em,itemsep=2pt]
\item 取 $\cH_-$ 为 $\cH$ 的 $\cO(\PP^{1}\setminus\{0\})$-子模，由
  \[
    \bigl\{(\hbar\alpha)^{k}\hbar^{-1}\hbar^{-(n+1)\alpha}\mid 0\le k\le n\bigr\}
  \]
  生成。
\item 取 $\Omega_0$ 为 $\cE\otimes_{\cO_{\tilde\cM}\{\hbar\}}\cO_{\tilde\cM}\{\hbar,\hbar^{-1}\}$ 的平坦截面，使得它在 $0\in\tilde\cM$ 处的值为 $[\Omega]$。
\end{itemize}

\begin{proposition}
(1) $\cH=\cE_0\oplus\cH_-$。

(2) $\cH_-$ 关于 (2.15) 中定义的辛形式 $\overline{\Omega}$ 是迷向的。

(3) $\operatorname{Gr}$ 保持 $\cH_-$，平坦截面 $\Omega_0\in\hbar\cH_-$ 表示 $\operatorname{Gr}_0$ 的一个特征向量，其特征值为 $-1$，而 $\cE$ 的相应截面 $s_0'=\tau(\Omega_0\otimes1)$ 在 $0\in\tilde\cM$ 的某邻域内给出微小通用性。
\end{proposition}

作为推论，由定理 2.26，我们在 $0\in\tilde\cM$ 的某邻域内得到具有平坦单位元与 Euler 向量场 $E$ 的 Frobenius 流形结构。

\para{证明。}$\cE_0$ 由 $\cR^\vee|_{\{0\}\times\C}$ 在 $\hbar=0$ 的邻域内的截面给出（用的是定义 2.37 中给出的 $\cR^\vee$ 跨 $\hbar=0$ 的延拓）。如引理 2.41 的证明中所观察到的，这种截面有一组基由 $[W_0^{i}\Omega]$（$0\le i\le n$）给出，因此 $\cE_0$ 作为 $\C\{\hbar\}$-模由这些类生成。现在由 (2.37) 可写
\begin{align*}
  \sum_{j=0}^{n}\alpha^{j}\int_{\Xi_j}W_0^{i}e^{W_0/\hbar}\Omega
  &=\left(-\hbar^{2}\frac{\partial}{\partial\hbar}\right)^{i}
  \left(\sum_{j=0}^{n}\alpha^{j}\int_{\Xi_j}e^{W_0/\hbar}\Omega\right)\\
  &=\left(-\hbar^{2}\frac{\partial}{\partial\hbar}\right)^{i}(\xi).
\end{align*}

% [原文 p.85]
由 $\xi$ 的显式公式 (2.36)，$(\hbar^{2}\partial_\hbar)^{i}\xi$ 的行为是
\[
  (\alpha\hbar)^{i}\hbar^{-(n+1)\alpha}
\]
再加上 $\cH_-$ 中的项。于是显然 $\cH=\cE_0\oplus\cH_-$，(1) 成立。

其次注意，由于 $\alpha^{i}$ 是 $\cR^\vee$ 的平坦截面，$((-)^{*}\alpha^{i},\alpha^{j})$ 是常数。此外，由于 $\hbar^{-(n+1)\alpha}\alpha^{i}$ 是 $\cR^\vee$ 的单值截面，
\[
  \bigl((-)^{*}\hbar^{-(n+1)\alpha}\alpha^{i},\hbar^{-(n+1)\alpha}\alpha^{j}\bigr)
\]
是单值的。另一方面，把 $\hbar^{-(n+1)\alpha}$ 按 $\log\hbar$ 的幂展开，就得到
$\bigl((-)^{*}\hbar^{-(n+1)\alpha}\alpha^{i},\hbar^{-(n+1)\alpha}\alpha^{j}\bigr)$ 的一个公式，其
$\hbar$-依赖性是 $\log\hbar$ 的一个多项式。于是，为了使它是单值的，它事实上必须与 $\hbar$ 无关。因此
\begin{align*}
  0&=\hbar\partial_\hbar\bigl((-)^{*}(\hbar^{-(n+1)\alpha}\alpha^{i}),\hbar^{-(n+1)\alpha}\alpha^{j}\bigr)\\
  &=\bigl(-(-)^{*}((n+1)\hbar^{-(n+1)\alpha}\alpha^{i+1}),\hbar^{-(n+1)\alpha}\alpha^{j}\bigr)\\
  &\quad+\bigl((-)^{*}(\hbar^{-(n+1)\alpha}\alpha^{i}),-(n+1)\hbar^{-(n+1)\alpha}\alpha^{j+1}\bigr).
\end{align*}
于是事实上
\begin{equation}
  \bigl((-)^{*}(\hbar^{-(n+1)\alpha}\alpha^{i+1}),\hbar^{-(n+1)\alpha}\alpha^{j}\bigr)
  =-\bigl((-)^{*}(\hbar^{-(n+1)\alpha}\alpha^{i}),\hbar^{-(n+1)\alpha}\alpha^{j+1}\bigr). \tag{2.39}
\end{equation}
由此推出
\begin{equation}
  (\hbar^{-(n+1)\alpha}\alpha^{i},\hbar^{-(n+1)\alpha}\alpha^{j})_\cE
  =\begin{cases}
    0 & i+j>n,\\
    (\text{Constant})\cdot\hbar^{-n} & i+j\le n.
  \end{cases} \tag{2.40}
\end{equation}
于是为证明 $\cH_-$ 是迷向的，只需注意对 $k,\ell\ge1$，
\begin{align*}
  \overline{\Omega}\bigl((\hbar\alpha)^{i}\hbar^{-(n+1)\alpha-k},
  (\hbar\alpha)^{j}\hbar^{-(n+1)\alpha-\ell}\bigr)
  &=\operatorname{Res}_{\hbar=0}\bigl((\hbar\alpha)^{i}\hbar^{-(n+1)\alpha-k},
  (\hbar\alpha)^{j}\hbar^{-(n+1)\alpha-\ell}\bigr)_\cE d\hbar\\
  &=\begin{cases}
    0 & i+j>n\\
    \operatorname{Res}_{\hbar=0}\text{Constant}\cdot\hbar^{i+j-k-\ell-n}d\hbar & i+j\le n
  \end{cases}\\
  &=0.
\end{align*}

至于 (3)，注意对 $m\ge1$，
\begin{align*}
  \operatorname{Gr}\bigl((\hbar\alpha)^{k}\hbar^{-(n+1)\alpha}\hbar^{-m}\bigr)
  &=\nabla^{GM}_{\hbar\partial_\hbar+E}\bigl(\hbar^{k-m}\alpha^{k}\hbar^{-(n+1)\alpha}\bigr)
  -(\hbar\alpha)^{k}\hbar^{-(n+1)\alpha}\hbar^{-m}\\
  &=\bigl(k-m-(n+1)\alpha-1\bigr)(\hbar\alpha)^{k}\hbar^{-(n+1)\alpha}\hbar^{-m}\\
  &\in\cH_-,
\end{align*}
故 $\operatorname{Gr}$ 保持 $\cH_-$。我们还看到 $\Omega_0\in\hbar\cH_-$，因为表示平坦截面 $\Omega_0$ 的 $\xi$ 形如 $\xi=\hbar^{-(n+1)\alpha}\bigl(1+O(\hbar^{-(n+1)d})\bigr)$。此外，
\begin{equation}
  \operatorname{Gr}_0\bigl((\hbar\alpha)^{k}\hbar^{-(n+1)\alpha}\bigr)
  =(k-1)(\hbar\alpha)^{k}\hbar^{-(n+1)\alpha}, \tag{2.41}
\end{equation}
故 $[\Omega_0]$ 是 $\operatorname{Gr}_0$ 的特征值为 $-1$ 的特征向量。

% [Tutor Note: 原文 p.85（本页）有两处可核对的问题。其一：“we need to show minversality” 中的 minversality 是 miniversality 的笔误（本页下文与 p.86、p.87 均作 miniversality），译文按正确拼写处理。其二：$\xi=\hbar^{-(n+1)\alpha}(1+O(\hbar^{-(n+1)d}))$ 中的 $d$ 疑为笔误：由 (2.36) 有 $\xi=\hbar^{-(n+1)\alpha}\sum_{d\ge0}\hbar^{-(n+1)d}\prod_{i=1}^{d}(\alpha+i)^{-(n+1)}$，故余项应为 $O(\hbar^{-(n+1)})$（$d=1$ 项即 $1/(\alpha+1)^{n+1}\hbar^{-(n+1)}$）。]

最后，我们需要证明微小通用性。特别地，令 $s_0'=\tau([\Omega_0]\otimes1)$。按定义，$s_0'$ 是这样得到的：把 $\cH_{\tilde\cM}$ 的常截面 $\Omega_0$（它在 $0\in\tilde\cM$ 处的值由 $[\Omega]$ 给出）写成和 $s_0'+s_0''$，其中 $s_0'$ 是 $\cE$ 的截面，$s_0''$ 是 $\cH_{-,\tilde\cM}$ 的截面。换言之，$s_0'$ 将由某个 $f\Omega$ 表示，其中 $f$ 是 $\check\cX\times\{|\hbar|<\epsilon\}$ 上的全纯函数，且在 $\check\cX$ 的纤维上正则。由于

% [原文 p.86]
$\Omega_0=\hbar^{-(n+1)\alpha}\mod\cH_-$，而我们要求 $[f\Omega]\equiv\Omega_0\mod\cH_-$，故对 $f\Omega$ 的要求是
\[
  [f\Omega]=\hbar^{-(n+1)\alpha}\sum_{i=0}^{n}\varphi_i(\mathbf{t},\hbar^{-1})(\alpha\hbar)^{i},
\]
其中
\[
  \varphi_i(\mathbf{t},\hbar^{-1})=\delta_{0,i}
  +\sum_{j=1}^{\infty}\varphi_{i,j}(\mathbf{t})\hbar^{-j}.
\]
这里我们记 $\mathbf{t}=(t_0,\ldots,t_n)$。此外，由于 $[\Omega]\in\cE_0$，事实上 $s_0'=[\Omega]$ 在 $0\in\tilde\cM$ 处成立，故可以假设 $f|_{\mathbf{t}=0}\equiv1$。这两个条件唯一地确定了 $f\Omega$ 所表示的类。

现在一旦给定 $s_0'=[f\Omega]$，我们就用 (2.30) 与 (2.34) 计算，以验证 $0$ 处的微小通用性：
\[
  \hbar\nabla_{\partial/\partial t_i}[f\Omega]
  =\begin{cases}
    \bigl[(\hbar(\partial f/\partial t_i)+fW_0^{i})\Omega\bigr] & i\ne1\\
    \bigl[\bigl(\hbar(\partial f/\partial t_1)+\frac{f}{n+1}(W_0+\sum_{i=2}^{n}it_iW_0^{i})\bigr)\Omega\bigr] & i=1
  \end{cases}
\]
于是在 $0\in\tilde\cM$ 处的 $\cE/\hbar\cE$ 中，若 $i\ne1$ 这就是 $[W_0^{i}\Omega]$，若 $i=1$ 则是 $\frac{1}{n+1}[W_0\Omega]$。我们知道 $[W_0^{i}\Omega]$（$0\le i\le n$）构成 $\cE/\hbar\cE$ 在 $0\in\tilde\cM$ 处的一组基，因为 $W_0^{i}$（$0\le i\le n$）构成 $W_0$ 的 Milnor 环的一组基。于是我们在 $0\in\tilde\cM$ 处有微小通用性，从而在 $0\in\tilde\cM$ 的某邻域内也有微小通用性。

\subsubsection{$\PP^{n}$ 的镜像对称。}现在我们终于可以陈述 $X=\PP^{n}$ 的镜像对称了。总结一下：我们现在有两组数据，即 A-模型数据与 B-模型数据。A-模型数据由 $\tilde\cM^{A}$ 上的半无限 Hodge 结构变分 $\cE^{A}$、$\nabla^{A}$、$(\cdot,\cdot)_{\cE^{A}}$，连同 $0\in\tilde\cM^{A}$ 处的对置子空间（opposite subspace）$\cH^{A}_-\subseteq\cH^{A}$ 以及元素 $\Omega^{A}_0\in\hbar\cH^{A}_-$ 组成。这些数据来自 $\PP^{n}$ 的量子上同调。B-模型数据由完全同类的数据组成，只是这次带上 $B$ 的标号，来自前两小节在 $\tilde\cM^{B}$ 上的构造。

特别地，有
\[
  \hbar\cH^{A}_-/\cH^{A}_-\cong H^{*}(X,\C)
  =\C[T_1]/(T_1^{n+1}),
\]
后一个等号是作为带有经典上积的环而言的；当然 $T_i=T_1^{i}$，其中 $T_1$ 是 $H^{2}(X,\C)$ 的生成元。又
\[
  \hbar\cH^{B}_-/\cH^{B}_-\cong\C[\alpha]/(\alpha^{n+1}),
\]
其中 $\alpha^{i}$ 与 $\hbar^{-(n+1)\alpha}(\hbar\alpha)^{i}$ 等同。于是有一个显然同构的选取
\begin{equation}
  \begin{gathered}
    \hbar\cH^{A}_-/\cH^{A}_-\cong\hbar\cH^{B}_-/\cH^{B}_-\\
    T_1^{i}\mapsto\hbar^{-(n+1)\alpha}(\hbar\alpha)^{i}.
  \end{gathered} \tag{2.42}
\end{equation}
由这些数据我们得到 Barannikov 的周期映射 (2.18) $\Psi^{A}$ 与 $\Psi^{B}$，从而得到芽的映射 (2.19)
\[
  \psi^{A}:(\tilde\cM^{A},0)\to\hbar\cH^{A}_-/\cH^{A}_-,
\]
以及
\[
  \psi^{B}:(\tilde\cM^{B},0)\to\hbar\cH^{B}_-/\cH^{B}_-,
\]

% [原文 p.87]
它们在 $0\in\tilde\cM$ 附近由微小通用性给出局部同构。利用 (2.42) 对 $\hbar\cH^{A}_-/\cH^{A}_-$ 与 $\hbar\cH^{B}_-/\cH^{B}_-$ 的等同，我们得到镜像映射
\[
  m:(\tilde\cM^{A},0)\to(\tilde\cM^{B},0).
\]

\begin{theorem}[$\PP^{n}$ 的镜像对称]
$m$ 把 A-模型与 B-模型的半无限 Hodge 结构变分等同起来。更具体地说，(2.42) 中给出的等同产生同构
\[
  \cE^{A}\xrightarrow{\tau_A^{-1}}
  (\hbar\cH^{A}_-/\cH^{A}_-)\otimes_\C\cO_{\tilde\cM}\{\hbar\}
  \cong(\hbar\cH^{B}_-/\cH^{B}_-)\otimes_\C\cO_{\tilde\cM}\{\hbar\}
  \xrightarrow{\tau_B}\cE^{B},
\]
它等同了联络、内积与分次结构。此外，对置子空间 $\cH^{A}_-$ 与 $\cH^{B}_-$，以及元素 $\Omega^{A}_0$ 与 $\Omega^{B}_0$，在这些同构下被等同。
\end{theorem}

这里不给出证明。它由 Barannikov 在 [4] 中证明，尽管并未完全用这些措辞。他证明了相应 Frobenius 流形的同构，而半无限 Hodge 结构变分的数据可由 Frobenius 流形结构恢复。证明中的基本想法是：$\PP^{n}$ 的量子上同调的 Frobenius 流形是半单的，从而完全由它在一点（比如 $0\in\tilde\cM^{A}$）处的行为决定。在 $\tilde\cM^{B}$ 的 $0$ 处可以找到同样的结构，因此这两个 Frobenius 流形在 $0$ 的某个邻域内重合。

% [Tutor Note: 定理 2.44 的证明被略去（见 p.87 第二段）：Barannikov [4] 的论证是 Frobenius 流形层面的同构加上“半单性决定局部结构”这一事实，再借 Frobenius 流形结构可恢复 $\cE,\nabla,(\cdot,\cdot)$（定理 2.26 的逆向）回到半无限 Hodge 结构变分。核对锚点：(2.18)/(2.19) 的周期映射 $\Psi^{A},\Psi^{B}$、(2.24) 的完美配对，以及命题 2.43 的微小通用性；三者合起来才能从 $m$ 得到本定理的全部内容。]

这个陈述相当抽象，而对第 5 章的目的而言，把它改写成一个远为具体的形式是方便的。

\begin{proposition}
定理 2.44 等价于如下陈述：

存在 $0\in\tilde\cM^{B}$ 的一个邻域 $U$，以及 $\cE^{B}$ 在 $U$ 上的一个截面 $s$，它由 $\pi^{-1}(U)\times\{|\hbar|<\epsilon\}$ 上的一个形式 $f\Omega$ 定义，其中 $f$ 是 $\pi$ 的纤维上代数的全纯函数，且 $f|_{\pi^{-1}(0)\times\C}\equiv1$，并满足以下条件：

(1) 若写
\[
  s(\mathbf{t},\hbar)=\sum_{i=0}^{n}\alpha^{i}\int_{\Xi_i}e^{W/\hbar}f\Omega,
\]
则
\[
  s(\mathbf{t},\hbar)=\hbar^{-(n+1)\alpha}\sum_{i=0}^{n}
  \varphi_i(\mathbf{t},\hbar^{-1})(\alpha\hbar)^{i},
\]
其中函数 $\varphi_i$ 满足
\[
  \varphi_i(\mathbf{t},\hbar^{-1})=\delta_{0,i}
  +\sum_{j=1}^{\infty}\varphi_{i,j}(\mathbf{t})\hbar^{-j},
  \qquad 0\le i\le n.
\]
若令
\[
  y_i(\mathbf{t})=\varphi_{i,1}(\mathbf{t}),\qquad 0\le i\le n,
\]
则 $y_0,\ldots,y_n$ 构成 $0$ 的邻域内 $\tilde\cM^{B}$ 的一个坐标系，它们被称为平坦坐标。

(2) 若把 Givental $J$-函数写成
\[
  J(y_0,\ldots,y_n,\hbar^{-1})
  =\sum_{i=0}^{n}J_i(y_0,\ldots,y_n,\hbar^{-1})T_i,
\]

% [原文 p.88]
则在 $\C$-向量空间 $\C[[y_0,\ldots,y_n,\hbar^{-1}]]$ 中
\begin{equation}
  J_i=\varphi_i. \tag{2.43}
\end{equation}

(3) 在由 $y_i\mapsto y_i$ 给出的映射 $\tilde\cM^{A}\to\tilde\cM^{B}$ 下，Euler 向量场 $E_A$ 与 $E_B$ 被等同。

(4)
\[
  \int_{\PP^{n}}T_i\cup T_j
  =\bigl(\hbar^{-(n+1)\alpha}(\hbar\alpha)^{i},
  \hbar^{-(n+1)\alpha}(\hbar\alpha)^{j}\bigr)_{\cE^{B}}\Big|_{\hbar=\infty}.
\]
\end{proposition}

\para{证明。}首先假设定理 2.44 成立。我们在命题 2.43(3) 的证明中已经看到，$s_0'=\tau(\Omega_0\otimes1)$ 实际上恰好由 (1) 中截面 $s$ 的描述给出。特别地，映射 $\psi_B:(\tilde\cM^{B},0)\to\hbar\cH^{B}_-/\cH^{B}_-$ 由
\[
  \psi_B(\mathbf{t})=\hbar\bigl(s(\mathbf{t},\hbar^{-1})-\hbar^{-(n+1)\alpha}\bigr)\mod\cH^{B}_-
\]
给出。于是用 $\hbar\cH^{B}_-/\cH^{B}_-$ 的基 $\{\hbar^{-(n+1)\alpha}(\hbar\alpha)^{i}\}$ 表示，我们看到
\[
  \psi_B(\mathbf{t})=\sum_{i=0}^{n}\varphi_{i,1}(\mathbf{t})
  \hbar^{-(n+1)\alpha}(\hbar\alpha)^{i}.
\]
这证明了 (1)。另一方面，由 (2.22)，
\[
  \psi_A(y_0,\ldots,y_n)=\sum_{i=0}^{n}y_iT_i.
\]
于是镜像映射 $m:\tilde\cM^{A}\to\tilde\cM^{B}$ 恰好由 $y_i\mapsto y_i$ 给出。由例 2.27，$\tau_A([T_0])=J_{\PP^{n}}$，而 $\tau_B([\Omega_0])=s$，于是得到 (2.43)。此外，(3) 是显然的，因为 Euler 向量场在镜像映射下重合。再者，由命题 2.24 与例 2.27，(4) 中等式左右两端分别是 $(\cdot,\cdot)_{\cE^{A}}$ 与 $(\cdot,\cdot)_{\cE^{B}}$ 在 $\hbar\cH^{A}_-/\cH^{A}_-$ 与 $\hbar\cH^{B}_-/\cH^{B}_-$ 上诱导的配对，故它们在等同 (2.42) 下必须一致。

反之，设给定一个 $f\Omega$，它给出满足所有给定性质的 $\cE^{B}$ 的截面。映射 $y_i\mapsto y_i$ 定义了镜像映射
\[
  m:(\tilde\cM^{A},0)\to(\tilde\cM^{B},0).
\]
如注 2.29 中所观察到的，$J_{\PP^{n}}$ 完全决定了 $\mathbb{J}$，从而决定了嵌入
\[
  \mathbb{J}:\cE^{A}\to\cH^{A}_{\tilde\cM},
\]
类似地，由 $f\Omega$ 定义的 $\cE^{B}$ 的截面 $s$ 的微小通用性，$s$ 决定了嵌入
\[
  \cE^{B}\to\cH^{B}_{\tilde\cM}.
\]
等式 (2.43) 于是保证：在用 (2.42) 等同
\[
  \begin{gathered}
    \cH^{A}_{\tilde\cM}=(\hbar\cH^{A}_-/\cH^{A}_-)\otimes_\C\cO_{\tilde\cM}\{\hbar,\hbar^{-1}\}\\
    \cH^{B}_{\tilde\cM}=(\hbar\cH^{B}_-/\cH^{B}_-)\otimes_\C\cO_{\tilde\cM}\{\hbar,\hbar^{-1}\}
  \end{gathered}
\]
之后，子丛 $\cE^{A}\subseteq\cH^{A}_{\tilde\cM}$ 与 $\cE^{B}\subseteq\cH^{B}_{\tilde\cM}$ 被等同。由条件 (4)，$(\cdot,\cdot)_{\cE^{A}}$ 与 $(\cdot,\cdot)_{\cE^{B}}$ 被等同。其次，$[\Omega^{A}_0]=T_0$、$[\Omega^{B}_0]=\hbar^{-(n+1)\alpha}$，且 $\cH^{A}_-$、$\cH^{B}_-$ 显然在 (2.42) 下被等同。最后，由 (3) 与命题 2.24(3)，

% [Tutor Note: 条件 (4) 右端在 $\hbar=\infty$ 处的取值需要说明：由 p.84 中 $\cH_-$ 的 $\cO(\PP^{1}\setminus\{0\})$-子模定义，$\hbar^{-(n+1)\alpha}(\hbar\alpha)^{i}$ 是 $\hbar\in\PP^{1}\setminus\{0\}$ 上的截面，故在 $\hbar=\infty$ 处有定义；而 $(\cdot,\cdot)_{\cE^{B}}$ 在 $\hbar=\infty$ 的行为由 (2.40) 给出（$i+j\le n$ 时为常数乘 $\hbar^{-n}$），故该极限存在，并与经典相交数 $\int_{\PP^{n}}T_i\cup T_j$ 一致（$i+j=n$ 时非零）。]

% [原文 p.89]
为证明 $\operatorname{Gr}^{A}$ 与 $\operatorname{Gr}^{B}$ 重合，只需证明 $\operatorname{Gr}^{A}_0$ 与 $\operatorname{Gr}^{B}_0$ 重合。但
\begin{align*}
  \operatorname{Gr}^{A}_0(T_i)
  &=\operatorname{Gr}^{A}(T_i)\mod\cH^{A}_-\\
  &=\operatorname{Gr}_E(T_i)\mod\cH^{A}_-\\
  &=\left(\frac{\deg T_i}{2}-1\right)T_i\\
  &=(i-1)T_i,
\end{align*}
而由 (2.41)，
\[
  \operatorname{Gr}^{B}_0\bigl(\hbar^{-(n+1)\alpha}(\alpha\hbar)^{i}\bigr)
  =(i-1)\hbar^{-(n+1)\alpha}(\alpha\hbar)^{i},
\]
故 $\operatorname{Gr}^{A}_0$ 与 $\operatorname{Gr}^{B}_0$ 重合。

\subsection{参考文献与延伸阅读}

关于 Gromov--Witten 不变量与量子上同调的讨论，很大程度上得益于 Fulton 与 Pandharipande 的综述 [28]。另一个有用的来源是 Cox 与 Katz 的书 [18]。关于 Frobenius 流形的大量信息，参见 Manin 的书 [76]。Coates--Iritani--Tseng [17] 与 Iritani [61] 的阐述，对于量子微分方程与半无限 Hodge 结构变分的讨论极为有益。

就 B-模型而言，我主要依据 Barannikov 的论文 [4]，并得到 Iritani [61] 与 Douai--Sabbah [20]、[21] 的帮助。关于 B-模型的一个重要的替代观点，见 Sabbah 的书 [103]。关于半无限 Hodge 结构变分中所涉及的某些结构的更一般讨论，见 Katzarkov、Kontsevich 与 Pantev 的文章 [67]。

% [Tutor Note: 本节只给出文献编号与评述，按本工程约定（BRIEF 第 5 条）不逐条翻译文献条目，完整书目由 Lead 在 c-refs 中以英文原貌给出。核对锚点：本节点名的编号为 [28]、[18]、[76]、[17]、[61]、[20]、[21]、[4]、[103]、[67]；而正文 p.78 引用过的 [102] 未在本节出现，需到原书参考文献（印刷 pp.307–312）中查。]

% [Tutor Note: 本章至此结束（印刷 p.89 正文后为空白）。本块末尾无图、无编号公式；下一块 c03a（tr-06）从 PDF p.106（印刷 p.91）的 §3.1 开始并写第 3 章标题。衔接检查：本块最后一句是证明的收尾（$\operatorname{Gr}^{A}_0=\operatorname{Gr}^{B}_0$），其后直接接 §2.3 的两段评述。]
